
\section{Background}



An \textit{agent} is typically defined as an autonomous entity that perceives its environment and acts upon it to achieve certain goals~\citep{Wooldridge2009}. A system composed of multiple interacting agents is referred to as a \textit{multi-agent system} (MAS). By coordinating their actions (whether cooperatively or competitively), agents in an MAS can solve problems that are beyond the capabilities of any individual agent or a centrally controlled system~\citep{Jennings1998}. The MAS paradigm, which originated in distributed artificial intelligence, emphasizes decentralized decision-making and has led to the development of various frameworks for agent communication, coordination, and negotiation over the past decades.

MAS techniques have been applied in a wide range of domains. For example, cooperative MAS deployments include fleets of warehouse robots coordinating storage and fulfillment~\citep{Wurman2008}, multi-robot teams in disaster response~\citep{Parker1998}, and LLM-based search and planning agents facilitating travel plans~\citep{choi2025atlas}. Competitive (adversarial) MAS arise when stakeholders have divergent objectives, e.g., defender–attacker resource allocation modeled as Stackelberg security games~\citep{Tambe2011}, or congestion management with selfish agents in routing games~\citep{Roughgarden2005}. Despite this breadth, many canonical MAS formalisms are computationally difficult to solve optimally even in realistic settings—for instance, planning in decentralized POMDPs is NEXP-complete even for finite horizons~\citep{bernstein2002complexity}, and computing Nash equilibria in general games is PPAD-complete~\citep{daskalakis2009complexity}. Amid this landscape, the \emph{distributed constraint optimization} (DCOP) framework provides a cooperative alternative that preserves key MAS characteristics—decentralized control, local objectives, and communication-based coordination—while admitting scalable exact/approximate algorithms and structure-exploiting methods for many practical problems~\citep{fioretto2018survey}. Consequently, we focus on DCOPs as the backbone for our study.

\paragraph{DCOP.} DCOPs involve a set of agents selecting local actions to maximize a global utility, typically formulated as the sum of local utility or constraint functions~\citep{fioretto2018survey}. Classical DCOP algorithms, such as complete search (e.g., ADOPT)~\citep{modi2005adopt}, message passing on pseudo-trees (e.g., DPOP)~\citep{petcu2005dpop}, and local/approximate methods (e.g., Max-Sum, DSA)~\citep{farinelli2008maxsum,zhang2005dsa}, trade off optimality guarantees, message and memory complexity, and anytime behavior, and they perform best when problem structure and communication protocols are explicitly defined and utilities are numeric and stationary. Building on this formalism, we extend the framework with LLM-based agents. Unlike the classical DCOP setting—where utilities are fixed symbolic functions and messages follow engineered schemas—our agents communicate in free-form natural language, and local constraints/objectives can be specified textually rather than solely as hand-coded numeric functions. This preserves the DCOP backbone while enabling systematic \emph{security} evaluations of LLM-specific vulnerabilities—e.g., prompt injection, communication poisoning.

\paragraph{Agent communication protocols.} Interacting agents in a MAS are required to communicate to achieve their objectives by utilizing a communication protocol that structures the rules and polices of communication between agents. Without this, performance degradation, inefficient token usage, and capability loss can emerge. Recent advancements in communication protocols such as the Agent2Agent (A2A) protocol~\citep{a2aprotocolProtocol} that assigns agent cards to specialized agents for more efficient collaboration initialization and allows agents to communicate with each other via messages, tools, and artifacts. Another established protocol is the Agent Communication Protocol (ACP)~\citep{githubGitHubIBMACP} which also enables agent-to-agent communication for low-latency communication. These protocols allow structured communication between agents and are vital to efficient and safe MAS. However, given that most communication protocols are developed by companies, we lack an open-source framework for testing and benchmarking communication protocols in a controllable and isolated environment which \sys satisfies.

\paragraph{Agent platforms and benchmarks.} LLM agents introduce significant security challenges such as indirect prompt injection attacks~\citep{greshake2023not} and context hijacking~\citep{bagdasarian2024airgapagent}. \citet{debenedetti2024agentdojo} proposed AgentDojo, which is a widely used dynamic environment to study agents' security and utility across domains such as Workspace, Travel, Slack, and Banking. Beyond single-agent benchmarks, \citet{abdelnabi2024cooperation} proposed a simulation benchmark for multi-agent interactive negotiation in a non-zero-sum game setup to study cooperation and competition dynamics. \citet{abdelnabi2025firewalls} studied security and privacy attacks in agent-to-agent communication where an AI assistant communicates with external parties to develop complex plans such as travel booking. Despite progress in this area, we lack a canonical platform that is easily extendable, configurable, and adaptable to study diverse multi-agent safety challenges, which is what we propose in our work. 

\paragraph{Multi-agent security.}
Increased adoption of LLM-based MAS has increased concerns about their security and privacy risks. Existing research demonstrates that known vulnerabilities such as prompt injection~\citep{greshake2023not} and jailbreaking~\citep{anil2024many} manifest more severely in multi-agent settings, where compromising a single agent enables the attack to propagate to all others~\citep{lee2024prompt}. Inspired by networking security challenges, recent work has also analyzed MAS protocols and introduced attack strategies including Agent-in-the-Middle~\citep{he2025red} and control-flow hijacking~\citep{triedman2025multi}. MASs have also been used to conduct attacks, such as jailbreaks, on other LLMs~\citep{rahman2025x,abdelnabi2025firewalls}.
% Attacks in multi-agent: safety, security, ... \eb{@Abhinav for you to finish}

\paragraph{Blackboards.} A \emph{blackboard} is a shared, structured workspace where heterogeneous agents post partial results, hypotheses, constraints, and goals for others to observe, refine, or refute. Historically, blackboard systems such as HEARSAY-II coordinated independent “knowledge sources” via a central store and scheduler, integrating evidence across abstraction layers to resolve uncertainty \citep{Erman1980,Nii1986a,Nii1986b}. In our multi-agent setting, multiple LLM-driven agents can similarly communicate and coordinate by appending proposals, commitments, and exception notes to a common log rather than engaging in bespoke pairwise messaging. Contemporary realizations span tuple-space designs (Linda-style generative communication) \citep{Gelernter1985}, append-only event logs for decoupled producers/consumers \citep{Kreps2011}, CRDT-backed shared documents for eventual consistency \citep{Shapiro2011}, and vector-indexed memories that enable retrieval-augmented reads \citep{Lewis2020}.

